% =============================================================================
\section{Energía en función del campo}\label{sec:campo}
% =============================================================================

Las expresiones anteriores asocian la energía a las cargas y la localizan donde están ellas. Igual que en el Tema~1 pasamos de la ley
de Coulomb, que habla de cargas, a la ley de Gauss, que habla de campos, vamos a expresar la energía solo en función del campo.

\subsection{El teorema de la energía del campo}

\begin{lema}[Dos herramientas]\label{lem:herramientas}
Para un campo escalar $\phi$ y un campo vectorial $\vA$ de clase $\mathcal C^1$:
\begin{enumerate}
  \item identidad del producto: $\divg(\phi\vA)=\grad\phi\cdot\vA+\phi\,\divg\vA$;
  \item teorema de la divergencia de Gauss: para una región acotada $V$ de frontera $S=\partial V$ con normal exterior,
        $\displaystyle\int_V\divg\vA\,\dd V=\oiint_S\vA\cdot\dS$.
\end{enumerate}
\end{lema}

\begin{teorema}[Energía electrostática en función del campo]\label{teo:energia-campo}
Para una distribución de carga localizada en el vacío,
\begin{equation}\label{eq:energia-campo}
  \boxed{\;U_e=\frac12\int_V\rho\,\phi\,\dd V=\frac{\eo}{2}\int_{\mathbb R^3}E^2\,\dd V\; ,}
\end{equation}
donde la segunda integral se extiende a todo el espacio.
\end{teorema}

\begin{proof}
\paso{Paso 1: eliminar $\rho$} Por la ley de Gauss diferencial, $\rho=\eo\,\divg\vE$, de modo que
$U_e=\dfrac{\eo}{2}\displaystyle\int_V\phi\,\divg\vE\,\dd V$. Como $\rho=0$ fuera de la distribución, podemos tomar como $V$ una bola $B_R$ de radio $R$
tan grande como queramos (\cref{fig:volumen}).

\paso{Paso 2: integrar por partes} Por el \cref{lem:herramientas}\,(i), $\phi\,\divg\vE=\divg(\phi\vE)-\grad\phi\cdot\vE=\divg(\phi\vE)+E^2$, ya que
$\grad\phi=-\vE$. Por tanto,
\[
  U_e=\frac{\eo}{2}\int_{B_R}\divg(\phi\vE)\,\dd V+\frac{\eo}{2}\int_{B_R}E^2\,\dd V .
\]
\paso{Paso 3: el teorema de la divergencia} Aplicándolo a la primera integral, con $S_R=\partial B_R$,
\[
  U_e=\frac{\eo}{2}\oiint_{S_R}\phi\,\vE\cdot\dS+\frac{\eo}{2}\int_{B_R}E^2\,\dd V .
\]
Esta igualdad vale para todo $R$ suficientemente grande, aunque cada término por separado dependa de $R$.

\paso{Paso 4: límite $R\to\infty$} Lejos de una distribución localizada, $\phi=\mathcal O(1/R)$ y $E=\mathcal O(1/R^2)$, mientras que el área de $S_R$ es $4\pi R^2$:
\[
  \abs{\oiint_{S_R}\phi\,\vE\cdot\dS}\leq\max_{S_R}\abs\phi\,\max_{S_R}\abs{\vE}\,4\pi R^2=\mathcal O\Big(\frac1R\cdot\frac{1}{R^2}\cdot R^2\Big)=\mathcal O\Big(\frac1R\Big)\xrightarrow{R\to\infty}0 .
\]
(Si la distribución es neutra, $\phi$ y $E$ decaen aún más deprisa.) La segunda integral converge, pues, a la integral de $E^2$ a todo el
espacio, y se obtiene \eqref{eq:energia-campo}.
\end{proof}

% Figura: paso de la energía en función de las cargas a la energía en función del
% campo. Volumen V con la carga, esfera S_R de radio R -> infinito y órdenes de magnitud.
\begin{figure}[t]
\centering
\begin{tikzpicture}[scale=1.0]
  % esfera grande S_R
  \draw[gauss] (0,0) circle (2.35);
  \node[etiq, text=acentoB] at (52:2.62) {$S_R$};
  \draw[texto!60!fondo, line width=0.4pt, -{Stealth[length=1.5mm]}] (0,0) -- (-35:2.35)
     node[pos=0.64, sloped, above, etiq, font=\footnotesize] {$R\to\infty$};
  % distribución de carga
  \fill[dielec] plot[smooth cycle, tension=0.85] coordinates
     {(-1.05,0.25) (-0.55,0.95) (0.35,0.85) (1.0,0.35) (0.85,-0.45) (0.05,-0.75) (-0.8,-0.5)};
  \draw[texto!70!fondo, line width=0.5pt] plot[smooth cycle, tension=0.85] coordinates
     {(-1.05,0.25) (-0.55,0.95) (0.35,0.85) (1.0,0.35) (0.85,-0.45) (0.05,-0.75) (-0.8,-0.5)};
  \node[etiq] at (-0.35,0.15) {$\rho(\vr)$};
  \node[etiq] at (0.55,-0.25) {$V$};
  \fill[cargapos] (0.05,0.5) circle (0.04);
  \node[etiqt, anchor=west] at (0.1,0.5) {$\dd V$};
  % elemento de superficie y campo lejano
  \coordinate (Q) at (125:2.35);
  \draw[vecn] (Q) -- ($(Q)+(125:0.75)$) node[anchor=south east, etiq] {$\dS$};
  \draw[vecE] (Q) -- ($(Q)+(95:0.65)$) node[anchor=south, etiq, text=campoE] {$\vE$};
  % tabla de escalas
  \node[anchor=west, align=left, font=\small, text=texto] at (3.05,0.35)
    {$\phi\sim\dfrac{1}{R}$,\qquad $E\sim\dfrac{1}{R^{2}}$,\qquad $\dd S\sim R^{2}$\\[3mm]
     $\displaystyle\Big|\oiint_{S_R}\phi\,\vE\cdot\dS\Big|\sim\frac1R\cdot\frac{1}{R^2}\cdot R^2=\frac1R\ \xrightarrow{R\to\infty}\ 0$};
\end{tikzpicture}
\caption{De las cargas a los campos. La energía $\frac12\int_V\rho\phi\,\dd V$ puede extenderse a cualquier volumen que contenga
a $V$, pues fuera de él $\rho=0$. Tomando una esfera $S_R$ de radio $R\to\infty$, el término de superficie que aparece al aplicar el
teorema de la divergencia decae como $1/R$ (para una distribución con carga neta; aún más deprisa si es neutra) y solo
sobrevive la integral de $E^2$ a todo el espacio (\cref{teo:energia-campo}).}
\label{fig:volumen}
\end{figure}


\begin{definicion}[Densidad de energía electrostática]\label{def:densidad}
Se llama \emph{densidad de energía electrostática} a
\begin{equation}\label{eq:densidad}
  u_e(\vr)\defeq\frac{\eo}{2}\,E^2(\vr)\qquad[\si{J/m^3}],\qquad U_e=\int_{\mathbb R^3}u_e\,\dd V .
\end{equation}
\end{definicion}

\begin{corolario}[Positividad]\label{cor:positiva}
La energía electrostática total de una distribución continua es positiva o nula, $U_e\geq0$, y solo se anula si $\vE\equiv\vcero$, es decir, si
no hay carga.
\end{corolario}

\begin{proof}
El integrando de \eqref{eq:energia-campo} es no negativo, y la integral de una función continua no negativa solo se anula si la
función es idénticamente nula.
\end{proof}

\begin{observacion}[¿Dónde está la energía?]\label{obs:localizacion}
Las dos expresiones de \eqref{eq:energia-campo} dan el mismo número, pero cuentan historias distintas: la primera localiza la energía
donde hay carga; la segunda, en todo punto donde hay campo, aunque allí no haya carga (\cref{fig:corteza}). En electrostática ambas
interpretaciones son equivalentes. La segunda, sin embargo, es la que sobrevive fuera de la electrostática: una onda
electromagnética transporta energía por regiones vacías, y $u_e=\frac{\eo}{2}E^2$ sigue siendo válida para campos variables en el
tiempo~\cite{Griffiths,Jackson}.
\end{observacion}

\begin{ejemplo}[Corteza esférica, de nuevo]\label{ej:corteza-campo}
Para la corteza del \cref{ej:corteza}, el campo es nulo dentro y vale $\vE=\frac{Q}{4\pi\eo r^2}\uv r$ fuera. Con $\dd V=4\pi r^2\dd r$ (simetría
esférica),
\[
  U_e=\frac{\eo}{2}\int_a^\infty\Big(\frac{Q}{4\pi\eo r^2}\Big)^2 4\pi r^2\,\dd r=\frac{Q^2}{8\pi\eo}\int_a^\infty\frac{\dd r}{r^2}
     =\frac{Q^2}{8\pi\eo}\Big[-\frac1r\Big]_a^\infty=\frac{Q^2}{8\pi\eo a},
\]
como antes. La energía almacenada más allá de un radio $r$ es $Q^2/(8\pi\eo r)$: la mitad está fuera de $2a$ (\cref{fig:corteza}).
\end{ejemplo}

% Figura: energía de una corteza esférica conductora de radio a y carga Q:
% densidad de energía u_e(r) y fracción de la energía almacenada fuera de r.
\begin{figure}[t]
\centering
\begin{subfigure}[b]{0.3\textwidth}\centering
\begin{tikzpicture}[scale=1.0]
  \shade[ball color=conductorluz] (0,0) circle (1.0);
  \draw[texto, line width=0.5pt] (0,0) circle (1.0);
  \foreach \t in {0,30,...,330} \node[cpos, font=\scriptsize\bfseries] at (\t:1.2) {$+$};
  \draw[texto, line width=0.4pt, -{Stealth[length=1.5mm]}] (0,0) -- (35:1.0);
  \node[font=\small, text=black] at (35:0.45) [above left=-2pt] {$a$};
  \node[etiq] at (0,-1.6) {$Q$, \ $\phi=\dfrac{Q}{4\pi\eo a}$};
\end{tikzpicture}
\caption{Corteza cargada}
\end{subfigure}\hfill
\begin{subfigure}[b]{0.33\textwidth}\centering
\begin{tikzpicture}
\begin{axis}[width=5.0cm, height=4.3cm, axis lines=left, xmin=0, xmax=4, ymin=0, ymax=1.15,
   xlabel={$r/a$}, ylabel={$u_e/u_e(a^{+})$}, xtick={0,1,2,3,4}, ytick={0,0.5,1}]
  \addplot[curvaE, domain=1:4, samples=80] {1/x^4};
  \addplot[curvaE] coordinates {(0,0) (1,0)};
  \draw[equipot] (axis cs:1,0) -- (axis cs:1,1.15);
  \node[etiqt, anchor=west] at (axis cs:1.4,0.55) {$\propto r^{-4}$};
\end{axis}
\end{tikzpicture}
\caption{Densidad de energía}
\end{subfigure}\hfill
\begin{subfigure}[b]{0.33\textwidth}\centering
\begin{tikzpicture}
\begin{axis}[width=5.0cm, height=4.3cm, axis lines=left, xmin=0, xmax=6, ymin=0, ymax=1.1,
   xlabel={$r/a$}, ylabel={$U_e(>r)/U_e$}, xtick={0,1,2,4,6}, ytick={0,0.5,1}, grid=major]
  \addplot[curvaD, domain=1:6, samples=80] {1/x};
  \addplot[curvaD] coordinates {(0,1) (1,1)};
  \addplot[only marks, mark=*, mark size=1.4pt, color=acentoC] coordinates {(2,0.5)};
  \node[etiqf, font=\scriptsize, anchor=south west] at (axis cs:2.15,0.55) {$\tfrac12$ en $r>2a$};
\end{axis}
\end{tikzpicture}
\caption{Energía fuera de $r$}
\end{subfigure}
\caption{\Cref{ej:corteza}. (a) Corteza esférica de radio $a$ con carga $Q$, equipotencial. (b) La densidad de energía
$u_e=\frac{\eo}{2}E^2$ es nula dentro y decae como $r^{-4}$ fuera. (c) La energía almacenada fuera de una esfera de radio $r>a$
es $U_e(>r)=Q^2/(8\pi\eo r)$: la mitad de la energía total está más allá de $r=2a$, lejos de la carga. Es la imagen que da
la \cref{obs:localizacion}: la energía «vive» en el campo.}
\label{fig:corteza}
\end{figure}


\subsection{La energía no es aditiva}

El campo verifica el principio de superposición; la energía, que es cuadrática en el campo, no.

\begin{proposicion}[Término cruzado]\label{prop:no-aditiva}
Si $\vE=\vE_1+\vE_2$ es el campo de dos distribuciones de carga, con energías propias $U_1$ y $U_2$,
\begin{equation}\label{eq:no-aditiva}
  U_e=\frac{\eo}{2}\int\big(E_1^2+E_2^2+2\,\vE_1\cdot\vE_2\big)\dd V=U_1+U_2+\eo\int_{\mathbb R^3}\vE_1\cdot\vE_2\,\dd V\neq U_1+U_2 .
\end{equation}
El término cruzado es la energía de interacción entre ambas distribuciones,
\begin{equation}\label{eq:cruzado}
  \Uint=\eo\int_{\mathbb R^3}\vE_1\cdot\vE_2\,\dd V=\int\rho_2\,\phi_1\,\dd V=\int\rho_1\,\phi_2\,\dd V .
\end{equation}
\end{proposicion}

\begin{proof}
La primera igualdad es el desarrollo de $(\vE_1+\vE_2)^2$. Para la segunda, con $\vE_1=-\grad\phi_1$ y la misma integración por partes
del \cref{teo:energia-campo},
\[
  \eo\int_{B_R}\vE_1\cdot\vE_2\,\dd V=-\eo\int_{B_R}\grad\phi_1\cdot\vE_2\,\dd V=-\eo\oiint_{S_R}\phi_1\vE_2\cdot\dS+\eo\int_{B_R}\phi_1\,\divg\vE_2\,\dd V ,
\]
y al hacer $R\to\infty$ el término de superficie se anula mientras que $\eo\,\divg\vE_2=\rho_2$. Intercambiando los papeles de 1 y 2 se obtiene la
otra forma.
\end{proof}

\begin{ejemplo}[Dos cargas puntuales]\label{ej:termino-cruzado}
Para dos cargas puntuales $q_1$ y $q_2$ separadas una distancia $d$, $U_1$ y $U_2$ son las energías propias, infinitas, de cada carga; el
término cruzado, en cambio, es finito y vale $\int\rho_2\phi_1\,\dd V=q_2\,\phi_1(\vr_2)=\dfrac{q_1q_2}{4\pi\eo d}$: la energía de interacción de la
\cref{prop:puntuales}. Para cargas iguales, $\vE_1\cdot\vE_2$ es negativo dentro de la circunferencia que tiene por diámetro el segmento que
las une (allí los campos forman un ángulo obtuso) y positivo fuera; el balance es positivo (\cref{fig:superposicion}).
\end{ejemplo}

% Figura (simulación): la energía no es aditiva. Densidad de energía de dos cargas
% iguales y opuestas, y término cruzado eps0 E1.E2.
% Mapas: scripts/tema5.py -> img/t5-ue-*.png ; líneas: datos/t5-dos-cargas-*.tex
\begin{figure}[t]
\centering
\pgfplotsset{mapa/.style={marco, width=5.6cm, xmin=-2.6, xmax=2.6, ymin=-1.8, ymax=1.8,
   enlargelimits=false, axis on top}}
\begin{subfigure}[b]{0.325\textwidth}\centering
\begin{tikzpicture}
\begin{axis}[mapa]
  \addplot graphics[xmin=-2.6, xmax=2.6, ymin=-1.8, ymax=1.8] {t5-ue-iguales-\modotema};
  %% Generado por scripts/tema5.py -- no editar a mano
\addplot[lineaE, flechaen=0.604] coordinates {(-1.111,0.069) (-1.146,0.090) (-1.180,0.111) (-1.214,0.132) (-1.249,0.153) (-1.283,0.174) (-1.318,0.194) (-1.353,0.215) (-1.387,0.235) (-1.422,0.256) (-1.457,0.276) (-1.492,0.296) (-1.527,0.316) (-1.562,0.336) (-1.597,0.355) (-1.632,0.375) (-1.668,0.395) (-1.703,0.414) (-1.738,0.433) (-1.774,0.453) (-1.809,0.472) (-1.845,0.491) (-1.880,0.510) (-1.916,0.529) (-1.951,0.547) (-1.987,0.566) (-2.023,0.585) (-2.059,0.603) (-2.095,0.622) (-2.130,0.640) (-2.166,0.658) (-2.202,0.676) (-2.238,0.694) (-2.274,0.712) (-2.310,0.730) (-2.346,0.748) (-2.383,0.766) (-2.419,0.784) (-2.455,0.802) (-2.491,0.819) (-2.527,0.837) (-2.564,0.854) (-2.600,0.872)};
\addplot[lineaE, flechaen=0.661] coordinates {(-1.073,0.113) (-1.095,0.146) (-1.117,0.180) (-1.139,0.213) (-1.161,0.246) (-1.184,0.280) (-1.207,0.312) (-1.229,0.345) (-1.253,0.378) (-1.276,0.411) (-1.299,0.443) (-1.323,0.475) (-1.347,0.508) (-1.371,0.539) (-1.395,0.571) (-1.420,0.603) (-1.444,0.634) (-1.469,0.666) (-1.494,0.697) (-1.520,0.728) (-1.545,0.759) (-1.571,0.790) (-1.597,0.820) (-1.623,0.851) (-1.649,0.881) (-1.675,0.911) (-1.701,0.941) (-1.728,0.971) (-1.755,1.001) (-1.781,1.031) (-1.808,1.061) (-1.835,1.090) (-1.863,1.119) (-1.890,1.149) (-1.917,1.178) (-1.945,1.207) (-1.972,1.236) (-2.000,1.265) (-2.028,1.294) (-2.056,1.322) (-2.084,1.351) (-2.112,1.380) (-2.140,1.408) (-2.168,1.437) (-2.196,1.465) (-2.225,1.493) (-2.253,1.521) (-2.282,1.549) (-2.310,1.578) (-2.339,1.606) (-2.368,1.633) (-2.396,1.661) (-2.425,1.689) (-2.454,1.717) (-2.483,1.745) (-2.512,1.772) (-2.541,1.800)};
\addplot[lineaE, flechaen=0.500] coordinates {(-1.031,0.126) (-1.040,0.166) (-1.050,0.206) (-1.061,0.245) (-1.071,0.285) (-1.082,0.324) (-1.093,0.364) (-1.105,0.403) (-1.116,0.442) (-1.128,0.481) (-1.141,0.520) (-1.154,0.559) (-1.167,0.597) (-1.181,0.636) (-1.195,0.674) (-1.209,0.713) (-1.224,0.751) (-1.239,0.789) (-1.255,0.827) (-1.271,0.864) (-1.287,0.902) (-1.303,0.939) (-1.320,0.976) (-1.337,1.014) (-1.355,1.050) (-1.373,1.087) (-1.391,1.124) (-1.409,1.161) (-1.428,1.197) (-1.446,1.233) (-1.465,1.269) (-1.485,1.306) (-1.504,1.341) (-1.524,1.377) (-1.544,1.413) (-1.564,1.449) (-1.584,1.484) (-1.604,1.520) (-1.625,1.555) (-1.646,1.590) (-1.667,1.625) (-1.688,1.660) (-1.709,1.695) (-1.730,1.730) (-1.751,1.765) (-1.773,1.800)};
\addplot[lineaE, flechaen=0.500] coordinates {(-0.990,0.131) (-0.988,0.171) (-0.985,0.212) (-0.983,0.253) (-0.981,0.293) (-0.980,0.334) (-0.978,0.374) (-0.978,0.415) (-0.977,0.456) (-0.978,0.496) (-0.978,0.537) (-0.980,0.577) (-0.982,0.618) (-0.984,0.659) (-0.987,0.699) (-0.991,0.740) (-0.995,0.780) (-1.000,0.820) (-1.005,0.861) (-1.011,0.901) (-1.018,0.941) (-1.025,0.981) (-1.032,1.021) (-1.040,1.061) (-1.049,1.100) (-1.057,1.140) (-1.067,1.180) (-1.077,1.219) (-1.087,1.258) (-1.097,1.298) (-1.108,1.337) (-1.119,1.376) (-1.131,1.415) (-1.142,1.454) (-1.155,1.493) (-1.167,1.531) (-1.179,1.570) (-1.192,1.608) (-1.205,1.647) (-1.219,1.685) (-1.232,1.724) (-1.246,1.762) (-1.260,1.800)};
\addplot[lineaE, flechaen=0.500] coordinates {(-0.950,0.121) (-0.935,0.159) (-0.919,0.197) (-0.905,0.235) (-0.890,0.273) (-0.876,0.312) (-0.862,0.350) (-0.849,0.389) (-0.836,0.428) (-0.824,0.467) (-0.813,0.506) (-0.803,0.545) (-0.793,0.585) (-0.784,0.625) (-0.776,0.665) (-0.768,0.705) (-0.762,0.745) (-0.756,0.786) (-0.752,0.826) (-0.748,0.867) (-0.745,0.908) (-0.743,0.948) (-0.741,0.989) (-0.741,1.030) (-0.741,1.071) (-0.742,1.112) (-0.743,1.153) (-0.745,1.193) (-0.748,1.234) (-0.751,1.275) (-0.755,1.316) (-0.759,1.356) (-0.763,1.397) (-0.768,1.437) (-0.773,1.478) (-0.779,1.518) (-0.785,1.559) (-0.791,1.599) (-0.798,1.639) (-0.805,1.680) (-0.812,1.720) (-0.819,1.760) (-0.827,1.800)};
\addplot[lineaE, flechaen=0.457] coordinates {(-0.909,0.095) (-0.882,0.124) (-0.854,0.154) (-0.827,0.183) (-0.800,0.213) (-0.773,0.243) (-0.746,0.274) (-0.720,0.305) (-0.695,0.336) (-0.670,0.368) (-0.645,0.400) (-0.622,0.433) (-0.599,0.466) (-0.578,0.500) (-0.557,0.535) (-0.538,0.571) (-0.520,0.607) (-0.503,0.643) (-0.488,0.681) (-0.474,0.719) (-0.462,0.757) (-0.450,0.796) (-0.440,0.835) (-0.432,0.874) (-0.424,0.914) (-0.417,0.954) (-0.412,0.994) (-0.407,1.034) (-0.403,1.074) (-0.400,1.114) (-0.398,1.155) (-0.396,1.195) (-0.394,1.235) (-0.394,1.276) (-0.393,1.316) (-0.393,1.356) (-0.394,1.397) (-0.395,1.437) (-0.396,1.478) (-0.397,1.518) (-0.399,1.558) (-0.401,1.599) (-0.403,1.639) (-0.406,1.679) (-0.408,1.719) (-0.411,1.760) (-0.414,1.800)};
\addplot[lineaE, flechaen=0.500] coordinates {(0.000,0.000) (0.000,0.041) (0.000,0.082) (0.000,0.123) (0.000,0.164) (0.000,0.205) (0.000,0.246) (0.000,0.286) (0.000,0.327) (0.000,0.368) (0.000,0.409) (0.000,0.450) (0.000,0.491) (0.000,0.532) (0.000,0.573) (0.000,0.614) (0.000,0.655) (0.000,0.696) (0.000,0.736) (0.000,0.777) (0.000,0.818) (0.000,0.859) (0.000,0.900) (0.000,0.941) (0.000,0.982) (0.000,1.023) (0.000,1.064) (0.000,1.105) (0.000,1.145) (0.000,1.186) (0.000,1.227) (0.000,1.268) (0.000,1.309) (0.000,1.350) (0.000,1.391) (0.000,1.432) (0.000,1.473) (0.000,1.514) (0.000,1.555) (0.000,1.595) (0.000,1.636) (0.000,1.677) (0.000,1.718) (0.000,1.759) (0.000,1.800)};
\addplot[lineaE, flechaen=0.457] coordinates {(0.909,0.095) (0.882,0.124) (0.854,0.154) (0.827,0.183) (0.800,0.213) (0.773,0.243) (0.746,0.274) (0.720,0.305) (0.695,0.336) (0.670,0.368) (0.645,0.400) (0.622,0.433) (0.599,0.466) (0.578,0.500) (0.557,0.535) (0.538,0.571) (0.520,0.607) (0.503,0.643) (0.488,0.681) (0.474,0.719) (0.462,0.757) (0.450,0.796) (0.440,0.835) (0.432,0.874) (0.424,0.914) (0.417,0.954) (0.412,0.994) (0.407,1.034) (0.403,1.074) (0.400,1.114) (0.398,1.155) (0.396,1.195) (0.394,1.235) (0.394,1.276) (0.393,1.316) (0.393,1.356) (0.394,1.397) (0.395,1.437) (0.396,1.478) (0.397,1.518) (0.399,1.558) (0.401,1.599) (0.403,1.639) (0.406,1.679) (0.408,1.719) (0.411,1.760) (0.414,1.800)};
\addplot[lineaE, flechaen=0.500] coordinates {(0.950,0.121) (0.935,0.159) (0.919,0.197) (0.905,0.235) (0.890,0.273) (0.876,0.312) (0.862,0.350) (0.849,0.389) (0.836,0.428) (0.824,0.467) (0.813,0.506) (0.803,0.545) (0.793,0.585) (0.784,0.625) (0.776,0.665) (0.768,0.705) (0.762,0.745) (0.756,0.786) (0.752,0.826) (0.748,0.867) (0.745,0.908) (0.743,0.948) (0.741,0.989) (0.741,1.030) (0.741,1.071) (0.742,1.112) (0.743,1.153) (0.745,1.193) (0.748,1.234) (0.751,1.275) (0.755,1.316) (0.759,1.356) (0.763,1.397) (0.768,1.437) (0.773,1.478) (0.779,1.518) (0.785,1.559) (0.791,1.599) (0.798,1.639) (0.805,1.680) (0.812,1.720) (0.819,1.760) (0.827,1.800)};
\addplot[lineaE, flechaen=0.500] coordinates {(0.990,0.131) (0.988,0.171) (0.985,0.212) (0.983,0.253) (0.981,0.293) (0.980,0.334) (0.978,0.374) (0.978,0.415) (0.977,0.456) (0.978,0.496) (0.978,0.537) (0.980,0.577) (0.982,0.618) (0.984,0.659) (0.987,0.699) (0.991,0.740) (0.995,0.780) (1.000,0.820) (1.005,0.861) (1.011,0.901) (1.018,0.941) (1.025,0.981) (1.032,1.021) (1.040,1.061) (1.049,1.100) (1.057,1.140) (1.067,1.180) (1.077,1.219) (1.087,1.258) (1.097,1.298) (1.108,1.337) (1.119,1.376) (1.131,1.415) (1.142,1.454) (1.155,1.493) (1.167,1.531) (1.179,1.570) (1.192,1.608) (1.205,1.647) (1.219,1.685) (1.232,1.724) (1.246,1.762) (1.260,1.800)};
\addplot[lineaE, flechaen=0.500] coordinates {(1.031,0.126) (1.040,0.166) (1.050,0.206) (1.061,0.245) (1.071,0.285) (1.082,0.324) (1.093,0.364) (1.105,0.403) (1.116,0.442) (1.128,0.481) (1.141,0.520) (1.154,0.559) (1.167,0.597) (1.181,0.636) (1.195,0.674) (1.209,0.713) (1.224,0.751) (1.239,0.789) (1.255,0.827) (1.271,0.864) (1.287,0.902) (1.303,0.939) (1.320,0.976) (1.337,1.014) (1.355,1.050) (1.373,1.087) (1.391,1.124) (1.409,1.161) (1.428,1.197) (1.446,1.233) (1.465,1.269) (1.485,1.306) (1.504,1.341) (1.524,1.377) (1.544,1.413) (1.564,1.449) (1.584,1.484) (1.604,1.520) (1.625,1.555) (1.646,1.590) (1.667,1.625) (1.688,1.660) (1.709,1.695) (1.730,1.730) (1.751,1.765) (1.773,1.800)};
\addplot[lineaE, flechaen=0.661] coordinates {(1.073,0.113) (1.095,0.146) (1.117,0.180) (1.139,0.213) (1.161,0.246) (1.184,0.280) (1.207,0.312) (1.229,0.345) (1.253,0.378) (1.276,0.411) (1.299,0.443) (1.323,0.475) (1.347,0.508) (1.371,0.539) (1.395,0.571) (1.420,0.603) (1.444,0.634) (1.469,0.666) (1.494,0.697) (1.520,0.728) (1.545,0.759) (1.571,0.790) (1.597,0.820) (1.623,0.851) (1.649,0.881) (1.675,0.911) (1.701,0.941) (1.728,0.971) (1.755,1.001) (1.781,1.031) (1.808,1.061) (1.835,1.090) (1.863,1.119) (1.890,1.149) (1.917,1.178) (1.945,1.207) (1.972,1.236) (2.000,1.265) (2.028,1.294) (2.056,1.322) (2.084,1.351) (2.112,1.380) (2.140,1.408) (2.168,1.437) (2.196,1.465) (2.225,1.493) (2.253,1.521) (2.282,1.549) (2.310,1.578) (2.339,1.606) (2.368,1.633) (2.396,1.661) (2.425,1.689) (2.454,1.717) (2.483,1.745) (2.512,1.772) (2.541,1.800)};
\addplot[lineaE, flechaen=0.604] coordinates {(1.111,0.069) (1.146,0.090) (1.180,0.111) (1.214,0.132) (1.249,0.153) (1.283,0.174) (1.318,0.194) (1.353,0.215) (1.387,0.235) (1.422,0.256) (1.457,0.276) (1.492,0.296) (1.527,0.316) (1.562,0.336) (1.597,0.355) (1.632,0.375) (1.668,0.395) (1.703,0.414) (1.738,0.433) (1.774,0.453) (1.809,0.472) (1.845,0.491) (1.880,0.510) (1.916,0.529) (1.951,0.547) (1.987,0.566) (2.023,0.585) (2.059,0.603) (2.095,0.622) (2.130,0.640) (2.166,0.658) (2.202,0.676) (2.238,0.694) (2.274,0.712) (2.310,0.730) (2.346,0.748) (2.383,0.766) (2.419,0.784) (2.455,0.802) (2.491,0.819) (2.527,0.837) (2.564,0.854) (2.600,0.872)};
\addplot[lineaE, flechaen=0.604] coordinates {(-1.111,-0.069) (-1.146,-0.090) (-1.180,-0.111) (-1.214,-0.132) (-1.249,-0.153) (-1.283,-0.174) (-1.318,-0.194) (-1.353,-0.215) (-1.387,-0.235) (-1.422,-0.256) (-1.457,-0.276) (-1.492,-0.296) (-1.527,-0.316) (-1.562,-0.336) (-1.597,-0.355) (-1.632,-0.375) (-1.668,-0.395) (-1.703,-0.414) (-1.738,-0.433) (-1.774,-0.453) (-1.809,-0.472) (-1.845,-0.491) (-1.880,-0.510) (-1.916,-0.529) (-1.951,-0.547) (-1.987,-0.566) (-2.023,-0.585) (-2.059,-0.603) (-2.095,-0.622) (-2.130,-0.640) (-2.166,-0.658) (-2.202,-0.676) (-2.238,-0.694) (-2.274,-0.712) (-2.310,-0.730) (-2.346,-0.748) (-2.383,-0.766) (-2.419,-0.784) (-2.455,-0.802) (-2.491,-0.819) (-2.527,-0.837) (-2.564,-0.854) (-2.600,-0.872)};
\addplot[lineaE, flechaen=0.661] coordinates {(-1.073,-0.113) (-1.095,-0.146) (-1.117,-0.180) (-1.139,-0.213) (-1.161,-0.246) (-1.184,-0.280) (-1.207,-0.312) (-1.229,-0.345) (-1.253,-0.378) (-1.276,-0.411) (-1.299,-0.443) (-1.323,-0.475) (-1.347,-0.508) (-1.371,-0.539) (-1.395,-0.571) (-1.420,-0.603) (-1.444,-0.634) (-1.469,-0.666) (-1.494,-0.697) (-1.520,-0.728) (-1.545,-0.759) (-1.571,-0.790) (-1.597,-0.820) (-1.623,-0.851) (-1.649,-0.881) (-1.675,-0.911) (-1.701,-0.941) (-1.728,-0.971) (-1.755,-1.001) (-1.781,-1.031) (-1.808,-1.061) (-1.835,-1.090) (-1.863,-1.119) (-1.890,-1.149) (-1.917,-1.178) (-1.945,-1.207) (-1.972,-1.236) (-2.000,-1.265) (-2.028,-1.294) (-2.056,-1.322) (-2.084,-1.351) (-2.112,-1.380) (-2.140,-1.408) (-2.168,-1.437) (-2.196,-1.465) (-2.225,-1.493) (-2.253,-1.521) (-2.282,-1.549) (-2.310,-1.578) (-2.339,-1.606) (-2.368,-1.633) (-2.396,-1.661) (-2.425,-1.689) (-2.454,-1.717) (-2.483,-1.745) (-2.512,-1.772) (-2.541,-1.800)};
\addplot[lineaE, flechaen=0.500] coordinates {(-1.031,-0.126) (-1.040,-0.166) (-1.050,-0.206) (-1.061,-0.245) (-1.071,-0.285) (-1.082,-0.324) (-1.093,-0.364) (-1.105,-0.403) (-1.116,-0.442) (-1.128,-0.481) (-1.141,-0.520) (-1.154,-0.559) (-1.167,-0.597) (-1.181,-0.636) (-1.195,-0.674) (-1.209,-0.713) (-1.224,-0.751) (-1.239,-0.789) (-1.255,-0.827) (-1.271,-0.864) (-1.287,-0.902) (-1.303,-0.939) (-1.320,-0.976) (-1.337,-1.014) (-1.355,-1.050) (-1.373,-1.087) (-1.391,-1.124) (-1.409,-1.161) (-1.428,-1.197) (-1.446,-1.233) (-1.465,-1.269) (-1.485,-1.306) (-1.504,-1.341) (-1.524,-1.377) (-1.544,-1.413) (-1.564,-1.449) (-1.584,-1.484) (-1.604,-1.520) (-1.625,-1.555) (-1.646,-1.590) (-1.667,-1.625) (-1.688,-1.660) (-1.709,-1.695) (-1.730,-1.730) (-1.751,-1.765) (-1.773,-1.800)};
\addplot[lineaE, flechaen=0.500] coordinates {(-0.990,-0.131) (-0.988,-0.171) (-0.985,-0.212) (-0.983,-0.253) (-0.981,-0.293) (-0.980,-0.334) (-0.978,-0.374) (-0.978,-0.415) (-0.977,-0.456) (-0.978,-0.496) (-0.978,-0.537) (-0.980,-0.577) (-0.982,-0.618) (-0.984,-0.659) (-0.987,-0.699) (-0.991,-0.740) (-0.995,-0.780) (-1.000,-0.820) (-1.005,-0.861) (-1.011,-0.901) (-1.018,-0.941) (-1.025,-0.981) (-1.032,-1.021) (-1.040,-1.061) (-1.049,-1.100) (-1.057,-1.140) (-1.067,-1.180) (-1.077,-1.219) (-1.087,-1.258) (-1.097,-1.298) (-1.108,-1.337) (-1.119,-1.376) (-1.131,-1.415) (-1.142,-1.454) (-1.155,-1.493) (-1.167,-1.531) (-1.179,-1.570) (-1.192,-1.608) (-1.205,-1.647) (-1.219,-1.685) (-1.232,-1.724) (-1.246,-1.762) (-1.260,-1.800)};
\addplot[lineaE, flechaen=0.500] coordinates {(-0.950,-0.121) (-0.935,-0.159) (-0.919,-0.197) (-0.905,-0.235) (-0.890,-0.273) (-0.876,-0.312) (-0.862,-0.350) (-0.849,-0.389) (-0.836,-0.428) (-0.824,-0.467) (-0.813,-0.506) (-0.803,-0.545) (-0.793,-0.585) (-0.784,-0.625) (-0.776,-0.665) (-0.768,-0.705) (-0.762,-0.745) (-0.756,-0.786) (-0.752,-0.826) (-0.748,-0.867) (-0.745,-0.908) (-0.743,-0.948) (-0.741,-0.989) (-0.741,-1.030) (-0.741,-1.071) (-0.742,-1.112) (-0.743,-1.153) (-0.745,-1.193) (-0.748,-1.234) (-0.751,-1.275) (-0.755,-1.316) (-0.759,-1.356) (-0.763,-1.397) (-0.768,-1.437) (-0.773,-1.478) (-0.779,-1.518) (-0.785,-1.559) (-0.791,-1.599) (-0.798,-1.639) (-0.805,-1.680) (-0.812,-1.720) (-0.819,-1.760) (-0.827,-1.800)};
\addplot[lineaE, flechaen=0.457] coordinates {(-0.909,-0.095) (-0.882,-0.124) (-0.854,-0.154) (-0.827,-0.183) (-0.800,-0.213) (-0.773,-0.243) (-0.746,-0.274) (-0.720,-0.305) (-0.695,-0.336) (-0.670,-0.368) (-0.645,-0.400) (-0.622,-0.433) (-0.599,-0.466) (-0.578,-0.500) (-0.557,-0.535) (-0.538,-0.571) (-0.520,-0.607) (-0.503,-0.643) (-0.488,-0.681) (-0.474,-0.719) (-0.462,-0.757) (-0.450,-0.796) (-0.440,-0.835) (-0.432,-0.874) (-0.424,-0.914) (-0.417,-0.954) (-0.412,-0.994) (-0.407,-1.034) (-0.403,-1.074) (-0.400,-1.114) (-0.398,-1.155) (-0.396,-1.195) (-0.394,-1.235) (-0.394,-1.276) (-0.393,-1.316) (-0.393,-1.356) (-0.394,-1.397) (-0.395,-1.437) (-0.396,-1.478) (-0.397,-1.518) (-0.399,-1.558) (-0.401,-1.599) (-0.403,-1.639) (-0.406,-1.679) (-0.408,-1.719) (-0.411,-1.760) (-0.414,-1.800)};
\addplot[lineaE, flechaen=0.500] coordinates {(0.000,-0.000) (0.000,-0.041) (0.000,-0.082) (0.000,-0.123) (0.000,-0.164) (0.000,-0.205) (0.000,-0.246) (0.000,-0.286) (0.000,-0.327) (0.000,-0.368) (0.000,-0.409) (0.000,-0.450) (0.000,-0.491) (0.000,-0.532) (0.000,-0.573) (0.000,-0.614) (0.000,-0.655) (0.000,-0.696) (0.000,-0.736) (0.000,-0.777) (0.000,-0.818) (0.000,-0.859) (0.000,-0.900) (0.000,-0.941) (0.000,-0.982) (0.000,-1.023) (0.000,-1.064) (0.000,-1.105) (0.000,-1.145) (0.000,-1.186) (0.000,-1.227) (0.000,-1.268) (0.000,-1.309) (0.000,-1.350) (0.000,-1.391) (0.000,-1.432) (0.000,-1.473) (0.000,-1.514) (0.000,-1.555) (0.000,-1.595) (0.000,-1.636) (0.000,-1.677) (0.000,-1.718) (0.000,-1.759) (0.000,-1.800)};
\addplot[lineaE, flechaen=0.457] coordinates {(0.909,-0.095) (0.882,-0.124) (0.854,-0.154) (0.827,-0.183) (0.800,-0.213) (0.773,-0.243) (0.746,-0.274) (0.720,-0.305) (0.695,-0.336) (0.670,-0.368) (0.645,-0.400) (0.622,-0.433) (0.599,-0.466) (0.578,-0.500) (0.557,-0.535) (0.538,-0.571) (0.520,-0.607) (0.503,-0.643) (0.488,-0.681) (0.474,-0.719) (0.462,-0.757) (0.450,-0.796) (0.440,-0.835) (0.432,-0.874) (0.424,-0.914) (0.417,-0.954) (0.412,-0.994) (0.407,-1.034) (0.403,-1.074) (0.400,-1.114) (0.398,-1.155) (0.396,-1.195) (0.394,-1.235) (0.394,-1.276) (0.393,-1.316) (0.393,-1.356) (0.394,-1.397) (0.395,-1.437) (0.396,-1.478) (0.397,-1.518) (0.399,-1.558) (0.401,-1.599) (0.403,-1.639) (0.406,-1.679) (0.408,-1.719) (0.411,-1.760) (0.414,-1.800)};
\addplot[lineaE, flechaen=0.500] coordinates {(0.950,-0.121) (0.935,-0.159) (0.919,-0.197) (0.905,-0.235) (0.890,-0.273) (0.876,-0.312) (0.862,-0.350) (0.849,-0.389) (0.836,-0.428) (0.824,-0.467) (0.813,-0.506) (0.803,-0.545) (0.793,-0.585) (0.784,-0.625) (0.776,-0.665) (0.768,-0.705) (0.762,-0.745) (0.756,-0.786) (0.752,-0.826) (0.748,-0.867) (0.745,-0.908) (0.743,-0.948) (0.741,-0.989) (0.741,-1.030) (0.741,-1.071) (0.742,-1.112) (0.743,-1.153) (0.745,-1.193) (0.748,-1.234) (0.751,-1.275) (0.755,-1.316) (0.759,-1.356) (0.763,-1.397) (0.768,-1.437) (0.773,-1.478) (0.779,-1.518) (0.785,-1.559) (0.791,-1.599) (0.798,-1.639) (0.805,-1.680) (0.812,-1.720) (0.819,-1.760) (0.827,-1.800)};
\addplot[lineaE, flechaen=0.500] coordinates {(0.990,-0.131) (0.988,-0.171) (0.985,-0.212) (0.983,-0.253) (0.981,-0.293) (0.980,-0.334) (0.978,-0.374) (0.978,-0.415) (0.977,-0.456) (0.978,-0.496) (0.978,-0.537) (0.980,-0.577) (0.982,-0.618) (0.984,-0.659) (0.987,-0.699) (0.991,-0.740) (0.995,-0.780) (1.000,-0.820) (1.005,-0.861) (1.011,-0.901) (1.018,-0.941) (1.025,-0.981) (1.032,-1.021) (1.040,-1.061) (1.049,-1.100) (1.057,-1.140) (1.067,-1.180) (1.077,-1.219) (1.087,-1.258) (1.097,-1.298) (1.108,-1.337) (1.119,-1.376) (1.131,-1.415) (1.142,-1.454) (1.155,-1.493) (1.167,-1.531) (1.179,-1.570) (1.192,-1.608) (1.205,-1.647) (1.219,-1.685) (1.232,-1.724) (1.246,-1.762) (1.260,-1.800)};
\addplot[lineaE, flechaen=0.500] coordinates {(1.031,-0.126) (1.040,-0.166) (1.050,-0.206) (1.061,-0.245) (1.071,-0.285) (1.082,-0.324) (1.093,-0.364) (1.105,-0.403) (1.116,-0.442) (1.128,-0.481) (1.141,-0.520) (1.154,-0.559) (1.167,-0.597) (1.181,-0.636) (1.195,-0.674) (1.209,-0.713) (1.224,-0.751) (1.239,-0.789) (1.255,-0.827) (1.271,-0.864) (1.287,-0.902) (1.303,-0.939) (1.320,-0.976) (1.337,-1.014) (1.355,-1.050) (1.373,-1.087) (1.391,-1.124) (1.409,-1.161) (1.428,-1.197) (1.446,-1.233) (1.465,-1.269) (1.485,-1.306) (1.504,-1.341) (1.524,-1.377) (1.544,-1.413) (1.564,-1.449) (1.584,-1.484) (1.604,-1.520) (1.625,-1.555) (1.646,-1.590) (1.667,-1.625) (1.688,-1.660) (1.709,-1.695) (1.730,-1.730) (1.751,-1.765) (1.773,-1.800)};
\addplot[lineaE, flechaen=0.661] coordinates {(1.073,-0.113) (1.095,-0.146) (1.117,-0.180) (1.139,-0.213) (1.161,-0.246) (1.184,-0.280) (1.207,-0.312) (1.229,-0.345) (1.253,-0.378) (1.276,-0.411) (1.299,-0.443) (1.323,-0.475) (1.347,-0.508) (1.371,-0.539) (1.395,-0.571) (1.420,-0.603) (1.444,-0.634) (1.469,-0.666) (1.494,-0.697) (1.520,-0.728) (1.545,-0.759) (1.571,-0.790) (1.597,-0.820) (1.623,-0.851) (1.649,-0.881) (1.675,-0.911) (1.701,-0.941) (1.728,-0.971) (1.755,-1.001) (1.781,-1.031) (1.808,-1.061) (1.835,-1.090) (1.863,-1.119) (1.890,-1.149) (1.917,-1.178) (1.945,-1.207) (1.972,-1.236) (2.000,-1.265) (2.028,-1.294) (2.056,-1.322) (2.084,-1.351) (2.112,-1.380) (2.140,-1.408) (2.168,-1.437) (2.196,-1.465) (2.225,-1.493) (2.253,-1.521) (2.282,-1.549) (2.310,-1.578) (2.339,-1.606) (2.368,-1.633) (2.396,-1.661) (2.425,-1.689) (2.454,-1.717) (2.483,-1.745) (2.512,-1.772) (2.541,-1.800)};
\addplot[lineaE, flechaen=0.604] coordinates {(1.111,-0.069) (1.146,-0.090) (1.180,-0.111) (1.214,-0.132) (1.249,-0.153) (1.283,-0.174) (1.318,-0.194) (1.353,-0.215) (1.387,-0.235) (1.422,-0.256) (1.457,-0.276) (1.492,-0.296) (1.527,-0.316) (1.562,-0.336) (1.597,-0.355) (1.632,-0.375) (1.668,-0.395) (1.703,-0.414) (1.738,-0.433) (1.774,-0.453) (1.809,-0.472) (1.845,-0.491) (1.880,-0.510) (1.916,-0.529) (1.951,-0.547) (1.987,-0.566) (2.023,-0.585) (2.059,-0.603) (2.095,-0.622) (2.130,-0.640) (2.166,-0.658) (2.202,-0.676) (2.238,-0.694) (2.274,-0.712) (2.310,-0.730) (2.346,-0.748) (2.383,-0.766) (2.419,-0.784) (2.455,-0.802) (2.491,-0.819) (2.527,-0.837) (2.564,-0.854) (2.600,-0.872)};
\addplot[lineaE, flechaen=0.592] coordinates {(-1.130,0.000) (-2.600,0.000)};
\addplot[lineaE, flechaen=0.592] coordinates {(1.130,0.000) (2.600,0.000)};

  \fill[cargapos] (axis cs:-1,0) circle[radius=0.075];
  \fill[cargapos] (axis cs:1,0) circle[radius=0.075];
\end{axis}
\end{tikzpicture}
\caption{$+q$ y $+q$: $u_e$}
\end{subfigure}\hfill
\begin{subfigure}[b]{0.325\textwidth}\centering
\begin{tikzpicture}
\begin{axis}[mapa]
  \addplot graphics[xmin=-2.6, xmax=2.6, ymin=-1.8, ymax=1.8] {t5-ue-opuestas-\modotema};
  %% Generado por scripts/tema5.py -- no editar a mano
\addplot[lineaE, flechaen=0.419] coordinates {(2.600,1.027) (2.568,1.002) (2.537,0.977) (2.505,0.952) (2.473,0.928) (2.441,0.903) (2.409,0.879) (2.376,0.855) (2.344,0.831) (2.311,0.807) (2.279,0.784) (2.246,0.760) (2.213,0.737) (2.180,0.714) (2.147,0.691) (2.114,0.668) (2.081,0.645) (2.048,0.622) (2.015,0.600) (1.981,0.578) (1.947,0.555) (1.914,0.533) (1.880,0.511) (1.846,0.490) (1.812,0.468) (1.778,0.447) (1.744,0.425) (1.710,0.404) (1.675,0.383) (1.641,0.362) (1.607,0.341) (1.572,0.321) (1.537,0.300) (1.503,0.280) (1.468,0.259) (1.433,0.239) (1.398,0.219) (1.363,0.199) (1.328,0.180) (1.293,0.160) (1.258,0.140) (1.222,0.121) (1.187,0.102) (1.152,0.082) (1.117,0.063)};
\addplot[lineaE, flechaen=0.581] coordinates {(-1.117,0.063) (-1.152,0.082) (-1.187,0.102) (-1.222,0.121) (-1.258,0.140) (-1.293,0.160) (-1.328,0.180) (-1.363,0.199) (-1.398,0.219) (-1.433,0.239) (-1.468,0.259) (-1.503,0.280) (-1.537,0.300) (-1.572,0.321) (-1.607,0.341) (-1.641,0.362) (-1.675,0.383) (-1.710,0.404) (-1.744,0.425) (-1.778,0.447) (-1.812,0.468) (-1.846,0.490) (-1.880,0.511) (-1.914,0.533) (-1.947,0.555) (-1.981,0.578) (-2.015,0.600) (-2.048,0.622) (-2.081,0.645) (-2.114,0.668) (-2.147,0.691) (-2.180,0.714) (-2.213,0.737) (-2.246,0.760) (-2.279,0.784) (-2.311,0.807) (-2.344,0.831) (-2.376,0.855) (-2.409,0.879) (-2.441,0.903) (-2.473,0.928) (-2.505,0.952) (-2.537,0.977) (-2.568,1.002) (-2.600,1.027)};
\addplot[lineaE, flechaen=0.159] coordinates {(2.112,1.800) (2.099,1.761) (2.086,1.723) (2.072,1.684) (2.058,1.646) (2.043,1.608) (2.028,1.570) (2.013,1.532) (1.997,1.495) (1.981,1.457) (1.964,1.420) (1.948,1.383) (1.930,1.346) (1.913,1.309) (1.895,1.272) (1.877,1.235) (1.858,1.199) (1.840,1.163) (1.820,1.127) (1.801,1.091) (1.781,1.055) (1.761,1.020) (1.741,0.984) (1.720,0.949) (1.699,0.914) (1.677,0.880) (1.656,0.845) (1.634,0.811) (1.611,0.777) (1.589,0.743) (1.566,0.709) (1.543,0.675) (1.519,0.642) (1.496,0.608) (1.472,0.575) (1.447,0.542) (1.423,0.510) (1.398,0.477) (1.373,0.445) (1.348,0.413) (1.323,0.381) (1.297,0.349) (1.272,0.317) (1.246,0.286) (1.220,0.254) (1.193,0.223) (1.167,0.192) (1.140,0.161) (1.114,0.130) (1.087,0.099)};
\addplot[lineaE, flechaen=0.841] coordinates {(-1.087,0.099) (-1.114,0.130) (-1.140,0.161) (-1.167,0.192) (-1.193,0.223) (-1.220,0.254) (-1.246,0.286) (-1.272,0.317) (-1.297,0.349) (-1.323,0.381) (-1.348,0.413) (-1.373,0.445) (-1.398,0.477) (-1.423,0.510) (-1.447,0.542) (-1.472,0.575) (-1.496,0.608) (-1.519,0.642) (-1.543,0.675) (-1.566,0.709) (-1.589,0.743) (-1.611,0.777) (-1.634,0.811) (-1.656,0.845) (-1.677,0.880) (-1.699,0.914) (-1.720,0.949) (-1.741,0.984) (-1.761,1.020) (-1.781,1.055) (-1.801,1.091) (-1.820,1.127) (-1.840,1.163) (-1.858,1.199) (-1.877,1.235) (-1.895,1.272) (-1.913,1.309) (-1.930,1.346) (-1.948,1.383) (-1.964,1.420) (-1.981,1.457) (-1.997,1.495) (-2.013,1.532) (-2.028,1.570) (-2.043,1.608) (-2.058,1.646) (-2.072,1.684) (-2.086,1.723) (-2.099,1.761) (-2.112,1.800)};
\addplot[lineaE, flechaen=0.500] coordinates {(1.459,1.800) (1.461,1.759) (1.463,1.719) (1.464,1.678) (1.465,1.637) (1.464,1.597) (1.463,1.556) (1.462,1.515) (1.460,1.474) (1.457,1.434) (1.454,1.393) (1.450,1.353) (1.445,1.312) (1.440,1.272) (1.434,1.232) (1.427,1.191) (1.420,1.151) (1.413,1.111) (1.405,1.071) (1.396,1.032) (1.387,0.992) (1.377,0.952) (1.367,0.913) (1.356,0.874) (1.345,0.835) (1.333,0.796) (1.321,0.757) (1.309,0.718) (1.296,0.679) (1.282,0.641) (1.268,0.603) (1.254,0.565) (1.239,0.527) (1.224,0.489) (1.209,0.451) (1.193,0.414) (1.177,0.376) (1.160,0.339) (1.144,0.302) (1.127,0.265) (1.110,0.228) (1.092,0.191) (1.075,0.154) (1.057,0.117)};
\addplot[lineaE, flechaen=0.500] coordinates {(-1.057,0.117) (-1.075,0.154) (-1.092,0.191) (-1.110,0.228) (-1.127,0.265) (-1.144,0.302) (-1.160,0.339) (-1.177,0.376) (-1.193,0.414) (-1.209,0.451) (-1.224,0.489) (-1.239,0.527) (-1.254,0.565) (-1.268,0.603) (-1.282,0.641) (-1.296,0.679) (-1.309,0.718) (-1.321,0.757) (-1.333,0.796) (-1.345,0.835) (-1.356,0.874) (-1.367,0.913) (-1.377,0.952) (-1.387,0.992) (-1.396,1.032) (-1.405,1.071) (-1.413,1.111) (-1.420,1.151) (-1.427,1.191) (-1.434,1.232) (-1.440,1.272) (-1.445,1.312) (-1.450,1.353) (-1.454,1.393) (-1.457,1.434) (-1.460,1.474) (-1.462,1.515) (-1.463,1.556) (-1.464,1.597) (-1.465,1.637) (-1.464,1.678) (-1.463,1.719) (-1.461,1.759) (-1.459,1.800)};
\addplot[lineaE, flechaen=0.500] coordinates {(0.904,1.800) (0.924,1.765) (0.943,1.728) (0.961,1.692) (0.979,1.655) (0.995,1.617) (1.010,1.579) (1.024,1.541) (1.037,1.502) (1.050,1.463) (1.061,1.424) (1.071,1.385) (1.081,1.345) (1.089,1.305) (1.097,1.265) (1.103,1.225) (1.109,1.184) (1.114,1.144) (1.119,1.103) (1.122,1.062) (1.125,1.022) (1.127,0.981) (1.128,0.940) (1.128,0.899) (1.128,0.858) (1.127,0.817) (1.125,0.777) (1.123,0.736) (1.120,0.695) (1.117,0.654) (1.112,0.614) (1.108,0.573) (1.103,0.533) (1.097,0.492) (1.091,0.452) (1.084,0.411) (1.078,0.371) (1.070,0.331) (1.063,0.291) (1.055,0.251) (1.046,0.211) (1.038,0.171) (1.029,0.131)};
\addplot[lineaE, flechaen=0.500] coordinates {(-1.029,0.131) (-1.038,0.171) (-1.046,0.211) (-1.055,0.251) (-1.063,0.291) (-1.070,0.331) (-1.078,0.371) (-1.084,0.411) (-1.091,0.452) (-1.097,0.492) (-1.103,0.533) (-1.108,0.573) (-1.112,0.614) (-1.117,0.654) (-1.120,0.695) (-1.123,0.736) (-1.125,0.777) (-1.127,0.817) (-1.128,0.858) (-1.128,0.899) (-1.128,0.940) (-1.127,0.981) (-1.125,1.022) (-1.122,1.062) (-1.119,1.103) (-1.114,1.144) (-1.109,1.184) (-1.103,1.225) (-1.097,1.265) (-1.089,1.305) (-1.081,1.345) (-1.071,1.385) (-1.061,1.424) (-1.050,1.463) (-1.037,1.502) (-1.024,1.541) (-1.010,1.579) (-0.995,1.617) (-0.979,1.655) (-0.961,1.692) (-0.943,1.728) (-0.924,1.765) (-0.904,1.800)};
\addplot[lineaE, flechaen=0.386, flechaen=0.614] coordinates {(-1.000,0.131) (-0.999,0.171) (-0.999,0.211) (-0.998,0.251) (-0.997,0.291) (-0.996,0.331) (-0.994,0.371) (-0.991,0.411) (-0.989,0.451) (-0.986,0.491) (-0.982,0.531) (-0.978,0.571) (-0.973,0.611) (-0.967,0.651) (-0.961,0.690) (-0.954,0.730) (-0.946,0.769) (-0.937,0.808) (-0.928,0.847) (-0.918,0.886) (-0.907,0.925) (-0.895,0.963) (-0.882,1.001) (-0.868,1.038) (-0.853,1.076) (-0.837,1.113) (-0.821,1.149) (-0.803,1.185) (-0.784,1.220) (-0.764,1.255) (-0.742,1.289) (-0.720,1.322) (-0.697,1.355) (-0.672,1.386) (-0.646,1.417) (-0.620,1.447) (-0.592,1.476) (-0.563,1.503) (-0.532,1.530) (-0.501,1.555) (-0.469,1.579) (-0.436,1.601) (-0.401,1.622) (-0.366,1.641) (-0.330,1.659) (-0.293,1.675) (-0.256,1.689) (-0.218,1.701) (-0.179,1.711) (-0.140,1.719) (-0.100,1.726) (-0.060,1.730) (-0.020,1.732) (0.020,1.732) (0.060,1.730) (0.100,1.726) (0.140,1.719) (0.179,1.711) (0.218,1.701) (0.256,1.689) (0.293,1.675) (0.330,1.659) (0.366,1.641) (0.401,1.622) (0.436,1.601) (0.469,1.579) (0.501,1.555) (0.532,1.530) (0.563,1.503) (0.592,1.476) (0.620,1.447) (0.646,1.417) (0.672,1.386) (0.697,1.355) (0.720,1.322) (0.742,1.289) (0.764,1.255) (0.784,1.220) (0.803,1.185) (0.821,1.149) (0.837,1.113) (0.853,1.076) (0.868,1.038) (0.882,1.001) (0.895,0.963) (0.907,0.925) (0.918,0.886) (0.928,0.847) (0.937,0.808) (0.946,0.769) (0.954,0.730) (0.961,0.690) (0.967,0.651) (0.973,0.611) (0.978,0.571) (0.982,0.531) (0.986,0.491) (0.989,0.451) (0.991,0.411) (0.994,0.371) (0.996,0.331) (0.997,0.291) (0.998,0.251) (0.999,0.211) (0.999,0.171) (1.000,0.131)};
\addplot[lineaE, flechaen=0.353, flechaen=0.647] coordinates {(-0.971,0.128) (-0.962,0.167) (-0.952,0.207) (-0.943,0.246) (-0.933,0.285) (-0.922,0.323) (-0.911,0.362) (-0.900,0.401) (-0.889,0.439) (-0.876,0.477) (-0.863,0.516) (-0.850,0.553) (-0.836,0.591) (-0.821,0.628) (-0.805,0.665) (-0.788,0.702) (-0.770,0.738) (-0.752,0.774) (-0.733,0.809) (-0.712,0.843) (-0.691,0.877) (-0.668,0.911) (-0.645,0.943) (-0.620,0.975) (-0.594,1.006) (-0.567,1.036) (-0.539,1.064) (-0.510,1.092) (-0.479,1.118) (-0.448,1.143) (-0.415,1.166) (-0.381,1.188) (-0.346,1.208) (-0.311,1.227) (-0.274,1.243) (-0.236,1.258) (-0.198,1.270) (-0.159,1.280) (-0.120,1.289) (-0.080,1.294) (-0.040,1.298) (-0.000,1.299) (0.040,1.298) (0.080,1.294) (0.120,1.289) (0.159,1.280) (0.198,1.270) (0.236,1.258) (0.274,1.243) (0.311,1.227) (0.346,1.208) (0.381,1.188) (0.415,1.166) (0.448,1.143) (0.479,1.118) (0.510,1.092) (0.539,1.064) (0.567,1.036) (0.594,1.006) (0.620,0.975) (0.645,0.943) (0.668,0.911) (0.691,0.877) (0.712,0.843) (0.733,0.809) (0.752,0.774) (0.770,0.738) (0.788,0.702) (0.805,0.665) (0.821,0.628) (0.836,0.591) (0.850,0.553) (0.863,0.516) (0.876,0.477) (0.889,0.439) (0.900,0.401) (0.911,0.362) (0.922,0.323) (0.933,0.285) (0.943,0.246) (0.952,0.207) (0.962,0.167) (0.971,0.128)};
\addplot[lineaE, flechaen=0.318, flechaen=0.682] coordinates {(-0.942,0.117) (-0.924,0.154) (-0.906,0.190) (-0.887,0.226) (-0.868,0.262) (-0.849,0.297) (-0.829,0.333) (-0.809,0.368) (-0.789,0.403) (-0.767,0.438) (-0.746,0.472) (-0.723,0.505) (-0.700,0.539) (-0.676,0.571) (-0.651,0.603) (-0.625,0.635) (-0.598,0.665) (-0.571,0.695) (-0.542,0.724) (-0.512,0.751) (-0.482,0.778) (-0.450,0.803) (-0.417,0.827) (-0.383,0.849) (-0.348,0.870) (-0.313,0.889) (-0.276,0.906) (-0.238,0.921) (-0.200,0.934) (-0.161,0.944) (-0.121,0.953) (-0.081,0.959) (-0.041,0.963) (-0.000,0.964) (0.041,0.963) (0.081,0.959) (0.121,0.953) (0.161,0.944) (0.200,0.934) (0.238,0.921) (0.276,0.906) (0.313,0.889) (0.348,0.870) (0.383,0.849) (0.417,0.827) (0.450,0.803) (0.482,0.778) (0.512,0.751) (0.542,0.724) (0.571,0.695) (0.598,0.665) (0.625,0.635) (0.651,0.603) (0.676,0.571) (0.700,0.539) (0.723,0.505) (0.746,0.472) (0.767,0.438) (0.789,0.403) (0.809,0.368) (0.829,0.333) (0.849,0.297) (0.868,0.262) (0.887,0.226) (0.906,0.190) (0.924,0.154) (0.942,0.117)};
\addplot[lineaE, flechaen=0.285, flechaen=0.715] coordinates {(-0.914,0.098) (-0.887,0.128) (-0.859,0.158) (-0.832,0.188) (-0.805,0.218) (-0.777,0.248) (-0.749,0.277) (-0.720,0.305) (-0.691,0.334) (-0.662,0.362) (-0.632,0.389) (-0.602,0.416) (-0.570,0.442) (-0.539,0.467) (-0.506,0.491) (-0.473,0.514) (-0.439,0.536) (-0.404,0.557) (-0.369,0.576) (-0.333,0.594) (-0.296,0.611) (-0.258,0.626) (-0.219,0.638) (-0.180,0.649) (-0.141,0.658) (-0.101,0.665) (-0.061,0.669) (-0.020,0.672) (0.020,0.672) (0.061,0.669) (0.101,0.665) (0.141,0.658) (0.180,0.649) (0.219,0.638) (0.258,0.626) (0.296,0.611) (0.333,0.594) (0.369,0.576) (0.404,0.557) (0.439,0.536) (0.473,0.514) (0.506,0.491) (0.539,0.467) (0.570,0.442) (0.602,0.416) (0.632,0.389) (0.662,0.362) (0.691,0.334) (0.720,0.305) (0.749,0.277) (0.777,0.248) (0.805,0.218) (0.832,0.188) (0.859,0.158) (0.887,0.128) (0.914,0.098)};
\addplot[lineaE, flechaen=0.255, flechaen=0.745] coordinates {(-0.882,0.063) (-0.847,0.082) (-0.812,0.101) (-0.776,0.119) (-0.741,0.138) (-0.705,0.156) (-0.669,0.174) (-0.633,0.191) (-0.597,0.208) (-0.560,0.225) (-0.523,0.241) (-0.486,0.257) (-0.449,0.271) (-0.412,0.285) (-0.374,0.298) (-0.335,0.311) (-0.297,0.322) (-0.258,0.332) (-0.219,0.340) (-0.180,0.348) (-0.140,0.354) (-0.100,0.358) (-0.060,0.361) (-0.020,0.363) (0.020,0.363) (0.060,0.361) (0.100,0.358) (0.140,0.354) (0.180,0.348) (0.219,0.340) (0.258,0.332) (0.297,0.322) (0.335,0.311) (0.374,0.298) (0.412,0.285) (0.449,0.271) (0.486,0.257) (0.523,0.241) (0.560,0.225) (0.597,0.208) (0.633,0.191) (0.669,0.174) (0.705,0.156) (0.741,0.138) (0.776,0.119) (0.812,0.101) (0.847,0.082) (0.882,0.063)};
\addplot[lineaE, flechaen=0.419] coordinates {(2.600,-1.027) (2.568,-1.002) (2.537,-0.977) (2.505,-0.952) (2.473,-0.928) (2.441,-0.903) (2.409,-0.879) (2.376,-0.855) (2.344,-0.831) (2.311,-0.807) (2.279,-0.784) (2.246,-0.760) (2.213,-0.737) (2.180,-0.714) (2.147,-0.691) (2.114,-0.668) (2.081,-0.645) (2.048,-0.622) (2.015,-0.600) (1.981,-0.578) (1.947,-0.555) (1.914,-0.533) (1.880,-0.511) (1.846,-0.490) (1.812,-0.468) (1.778,-0.447) (1.744,-0.425) (1.710,-0.404) (1.675,-0.383) (1.641,-0.362) (1.607,-0.341) (1.572,-0.321) (1.537,-0.300) (1.503,-0.280) (1.468,-0.259) (1.433,-0.239) (1.398,-0.219) (1.363,-0.199) (1.328,-0.180) (1.293,-0.160) (1.258,-0.140) (1.222,-0.121) (1.187,-0.102) (1.152,-0.082) (1.117,-0.063)};
\addplot[lineaE, flechaen=0.581] coordinates {(-1.117,-0.063) (-1.152,-0.082) (-1.187,-0.102) (-1.222,-0.121) (-1.258,-0.140) (-1.293,-0.160) (-1.328,-0.180) (-1.363,-0.199) (-1.398,-0.219) (-1.433,-0.239) (-1.468,-0.259) (-1.503,-0.280) (-1.537,-0.300) (-1.572,-0.321) (-1.607,-0.341) (-1.641,-0.362) (-1.675,-0.383) (-1.710,-0.404) (-1.744,-0.425) (-1.778,-0.447) (-1.812,-0.468) (-1.846,-0.490) (-1.880,-0.511) (-1.914,-0.533) (-1.947,-0.555) (-1.981,-0.578) (-2.015,-0.600) (-2.048,-0.622) (-2.081,-0.645) (-2.114,-0.668) (-2.147,-0.691) (-2.180,-0.714) (-2.213,-0.737) (-2.246,-0.760) (-2.279,-0.784) (-2.311,-0.807) (-2.344,-0.831) (-2.376,-0.855) (-2.409,-0.879) (-2.441,-0.903) (-2.473,-0.928) (-2.505,-0.952) (-2.537,-0.977) (-2.568,-1.002) (-2.600,-1.027)};
\addplot[lineaE, flechaen=0.159] coordinates {(2.112,-1.800) (2.099,-1.761) (2.086,-1.723) (2.072,-1.684) (2.058,-1.646) (2.043,-1.608) (2.028,-1.570) (2.013,-1.532) (1.997,-1.495) (1.981,-1.457) (1.964,-1.420) (1.948,-1.383) (1.930,-1.346) (1.913,-1.309) (1.895,-1.272) (1.877,-1.235) (1.858,-1.199) (1.840,-1.163) (1.820,-1.127) (1.801,-1.091) (1.781,-1.055) (1.761,-1.020) (1.741,-0.984) (1.720,-0.949) (1.699,-0.914) (1.677,-0.880) (1.656,-0.845) (1.634,-0.811) (1.611,-0.777) (1.589,-0.743) (1.566,-0.709) (1.543,-0.675) (1.519,-0.642) (1.496,-0.608) (1.472,-0.575) (1.447,-0.542) (1.423,-0.510) (1.398,-0.477) (1.373,-0.445) (1.348,-0.413) (1.323,-0.381) (1.297,-0.349) (1.272,-0.317) (1.246,-0.286) (1.220,-0.254) (1.193,-0.223) (1.167,-0.192) (1.140,-0.161) (1.114,-0.130) (1.087,-0.099)};
\addplot[lineaE, flechaen=0.841] coordinates {(-1.087,-0.099) (-1.114,-0.130) (-1.140,-0.161) (-1.167,-0.192) (-1.193,-0.223) (-1.220,-0.254) (-1.246,-0.286) (-1.272,-0.317) (-1.297,-0.349) (-1.323,-0.381) (-1.348,-0.413) (-1.373,-0.445) (-1.398,-0.477) (-1.423,-0.510) (-1.447,-0.542) (-1.472,-0.575) (-1.496,-0.608) (-1.519,-0.642) (-1.543,-0.675) (-1.566,-0.709) (-1.589,-0.743) (-1.611,-0.777) (-1.634,-0.811) (-1.656,-0.845) (-1.677,-0.880) (-1.699,-0.914) (-1.720,-0.949) (-1.741,-0.984) (-1.761,-1.020) (-1.781,-1.055) (-1.801,-1.091) (-1.820,-1.127) (-1.840,-1.163) (-1.858,-1.199) (-1.877,-1.235) (-1.895,-1.272) (-1.913,-1.309) (-1.930,-1.346) (-1.948,-1.383) (-1.964,-1.420) (-1.981,-1.457) (-1.997,-1.495) (-2.013,-1.532) (-2.028,-1.570) (-2.043,-1.608) (-2.058,-1.646) (-2.072,-1.684) (-2.086,-1.723) (-2.099,-1.761) (-2.112,-1.800)};
\addplot[lineaE, flechaen=0.500] coordinates {(1.459,-1.800) (1.461,-1.759) (1.463,-1.719) (1.464,-1.678) (1.465,-1.637) (1.464,-1.597) (1.463,-1.556) (1.462,-1.515) (1.460,-1.474) (1.457,-1.434) (1.454,-1.393) (1.450,-1.353) (1.445,-1.312) (1.440,-1.272) (1.434,-1.232) (1.427,-1.191) (1.420,-1.151) (1.413,-1.111) (1.405,-1.071) (1.396,-1.032) (1.387,-0.992) (1.377,-0.952) (1.367,-0.913) (1.356,-0.874) (1.345,-0.835) (1.333,-0.796) (1.321,-0.757) (1.309,-0.718) (1.296,-0.679) (1.282,-0.641) (1.268,-0.603) (1.254,-0.565) (1.239,-0.527) (1.224,-0.489) (1.209,-0.451) (1.193,-0.414) (1.177,-0.376) (1.160,-0.339) (1.144,-0.302) (1.127,-0.265) (1.110,-0.228) (1.092,-0.191) (1.075,-0.154) (1.057,-0.117)};
\addplot[lineaE, flechaen=0.500] coordinates {(-1.057,-0.117) (-1.075,-0.154) (-1.092,-0.191) (-1.110,-0.228) (-1.127,-0.265) (-1.144,-0.302) (-1.160,-0.339) (-1.177,-0.376) (-1.193,-0.414) (-1.209,-0.451) (-1.224,-0.489) (-1.239,-0.527) (-1.254,-0.565) (-1.268,-0.603) (-1.282,-0.641) (-1.296,-0.679) (-1.309,-0.718) (-1.321,-0.757) (-1.333,-0.796) (-1.345,-0.835) (-1.356,-0.874) (-1.367,-0.913) (-1.377,-0.952) (-1.387,-0.992) (-1.396,-1.032) (-1.405,-1.071) (-1.413,-1.111) (-1.420,-1.151) (-1.427,-1.191) (-1.434,-1.232) (-1.440,-1.272) (-1.445,-1.312) (-1.450,-1.353) (-1.454,-1.393) (-1.457,-1.434) (-1.460,-1.474) (-1.462,-1.515) (-1.463,-1.556) (-1.464,-1.597) (-1.465,-1.637) (-1.464,-1.678) (-1.463,-1.719) (-1.461,-1.759) (-1.459,-1.800)};
\addplot[lineaE, flechaen=0.500] coordinates {(0.904,-1.800) (0.924,-1.765) (0.943,-1.728) (0.961,-1.692) (0.979,-1.655) (0.995,-1.617) (1.010,-1.579) (1.024,-1.541) (1.037,-1.502) (1.050,-1.463) (1.061,-1.424) (1.071,-1.385) (1.081,-1.345) (1.089,-1.305) (1.097,-1.265) (1.103,-1.225) (1.109,-1.184) (1.114,-1.144) (1.119,-1.103) (1.122,-1.062) (1.125,-1.022) (1.127,-0.981) (1.128,-0.940) (1.128,-0.899) (1.128,-0.858) (1.127,-0.817) (1.125,-0.777) (1.123,-0.736) (1.120,-0.695) (1.117,-0.654) (1.112,-0.614) (1.108,-0.573) (1.103,-0.533) (1.097,-0.492) (1.091,-0.452) (1.084,-0.411) (1.078,-0.371) (1.070,-0.331) (1.063,-0.291) (1.055,-0.251) (1.046,-0.211) (1.038,-0.171) (1.029,-0.131)};
\addplot[lineaE, flechaen=0.500] coordinates {(-1.029,-0.131) (-1.038,-0.171) (-1.046,-0.211) (-1.055,-0.251) (-1.063,-0.291) (-1.070,-0.331) (-1.078,-0.371) (-1.084,-0.411) (-1.091,-0.452) (-1.097,-0.492) (-1.103,-0.533) (-1.108,-0.573) (-1.112,-0.614) (-1.117,-0.654) (-1.120,-0.695) (-1.123,-0.736) (-1.125,-0.777) (-1.127,-0.817) (-1.128,-0.858) (-1.128,-0.899) (-1.128,-0.940) (-1.127,-0.981) (-1.125,-1.022) (-1.122,-1.062) (-1.119,-1.103) (-1.114,-1.144) (-1.109,-1.184) (-1.103,-1.225) (-1.097,-1.265) (-1.089,-1.305) (-1.081,-1.345) (-1.071,-1.385) (-1.061,-1.424) (-1.050,-1.463) (-1.037,-1.502) (-1.024,-1.541) (-1.010,-1.579) (-0.995,-1.617) (-0.979,-1.655) (-0.961,-1.692) (-0.943,-1.728) (-0.924,-1.765) (-0.904,-1.800)};
\addplot[lineaE, flechaen=0.386, flechaen=0.614] coordinates {(-1.000,-0.131) (-0.999,-0.171) (-0.999,-0.211) (-0.998,-0.251) (-0.997,-0.291) (-0.996,-0.331) (-0.994,-0.371) (-0.991,-0.411) (-0.989,-0.451) (-0.986,-0.491) (-0.982,-0.531) (-0.978,-0.571) (-0.973,-0.611) (-0.967,-0.651) (-0.961,-0.690) (-0.954,-0.730) (-0.946,-0.769) (-0.937,-0.808) (-0.928,-0.847) (-0.918,-0.886) (-0.907,-0.925) (-0.895,-0.963) (-0.882,-1.001) (-0.868,-1.038) (-0.853,-1.076) (-0.837,-1.113) (-0.821,-1.149) (-0.803,-1.185) (-0.784,-1.220) (-0.764,-1.255) (-0.742,-1.289) (-0.720,-1.322) (-0.697,-1.355) (-0.672,-1.386) (-0.646,-1.417) (-0.620,-1.447) (-0.592,-1.476) (-0.563,-1.503) (-0.532,-1.530) (-0.501,-1.555) (-0.469,-1.579) (-0.436,-1.601) (-0.401,-1.622) (-0.366,-1.641) (-0.330,-1.659) (-0.293,-1.675) (-0.256,-1.689) (-0.218,-1.701) (-0.179,-1.711) (-0.140,-1.719) (-0.100,-1.726) (-0.060,-1.730) (-0.020,-1.732) (0.020,-1.732) (0.060,-1.730) (0.100,-1.726) (0.140,-1.719) (0.179,-1.711) (0.218,-1.701) (0.256,-1.689) (0.293,-1.675) (0.330,-1.659) (0.366,-1.641) (0.401,-1.622) (0.436,-1.601) (0.469,-1.579) (0.501,-1.555) (0.532,-1.530) (0.563,-1.503) (0.592,-1.476) (0.620,-1.447) (0.646,-1.417) (0.672,-1.386) (0.697,-1.355) (0.720,-1.322) (0.742,-1.289) (0.764,-1.255) (0.784,-1.220) (0.803,-1.185) (0.821,-1.149) (0.837,-1.113) (0.853,-1.076) (0.868,-1.038) (0.882,-1.001) (0.895,-0.963) (0.907,-0.925) (0.918,-0.886) (0.928,-0.847) (0.937,-0.808) (0.946,-0.769) (0.954,-0.730) (0.961,-0.690) (0.967,-0.651) (0.973,-0.611) (0.978,-0.571) (0.982,-0.531) (0.986,-0.491) (0.989,-0.451) (0.991,-0.411) (0.994,-0.371) (0.996,-0.331) (0.997,-0.291) (0.998,-0.251) (0.999,-0.211) (0.999,-0.171) (1.000,-0.131)};
\addplot[lineaE, flechaen=0.353, flechaen=0.647] coordinates {(-0.971,-0.128) (-0.962,-0.167) (-0.952,-0.207) (-0.943,-0.246) (-0.933,-0.285) (-0.922,-0.323) (-0.911,-0.362) (-0.900,-0.401) (-0.889,-0.439) (-0.876,-0.477) (-0.863,-0.516) (-0.850,-0.553) (-0.836,-0.591) (-0.821,-0.628) (-0.805,-0.665) (-0.788,-0.702) (-0.770,-0.738) (-0.752,-0.774) (-0.733,-0.809) (-0.712,-0.843) (-0.691,-0.877) (-0.668,-0.911) (-0.645,-0.943) (-0.620,-0.975) (-0.594,-1.006) (-0.567,-1.036) (-0.539,-1.064) (-0.510,-1.092) (-0.479,-1.118) (-0.448,-1.143) (-0.415,-1.166) (-0.381,-1.188) (-0.346,-1.208) (-0.311,-1.227) (-0.274,-1.243) (-0.236,-1.258) (-0.198,-1.270) (-0.159,-1.280) (-0.120,-1.289) (-0.080,-1.294) (-0.040,-1.298) (-0.000,-1.299) (0.040,-1.298) (0.080,-1.294) (0.120,-1.289) (0.159,-1.280) (0.198,-1.270) (0.236,-1.258) (0.274,-1.243) (0.311,-1.227) (0.346,-1.208) (0.381,-1.188) (0.415,-1.166) (0.448,-1.143) (0.479,-1.118) (0.510,-1.092) (0.539,-1.064) (0.567,-1.036) (0.594,-1.006) (0.620,-0.975) (0.645,-0.943) (0.668,-0.911) (0.691,-0.877) (0.712,-0.843) (0.733,-0.809) (0.752,-0.774) (0.770,-0.738) (0.788,-0.702) (0.805,-0.665) (0.821,-0.628) (0.836,-0.591) (0.850,-0.553) (0.863,-0.516) (0.876,-0.477) (0.889,-0.439) (0.900,-0.401) (0.911,-0.362) (0.922,-0.323) (0.933,-0.285) (0.943,-0.246) (0.952,-0.207) (0.962,-0.167) (0.971,-0.128)};
\addplot[lineaE, flechaen=0.318, flechaen=0.682] coordinates {(-0.942,-0.117) (-0.924,-0.154) (-0.906,-0.190) (-0.887,-0.226) (-0.868,-0.262) (-0.849,-0.297) (-0.829,-0.333) (-0.809,-0.368) (-0.789,-0.403) (-0.767,-0.438) (-0.746,-0.472) (-0.723,-0.505) (-0.700,-0.539) (-0.676,-0.571) (-0.651,-0.603) (-0.625,-0.635) (-0.598,-0.665) (-0.571,-0.695) (-0.542,-0.724) (-0.512,-0.751) (-0.482,-0.778) (-0.450,-0.803) (-0.417,-0.827) (-0.383,-0.849) (-0.348,-0.870) (-0.313,-0.889) (-0.276,-0.906) (-0.238,-0.921) (-0.200,-0.934) (-0.161,-0.944) (-0.121,-0.953) (-0.081,-0.959) (-0.041,-0.963) (-0.000,-0.964) (0.041,-0.963) (0.081,-0.959) (0.121,-0.953) (0.161,-0.944) (0.200,-0.934) (0.238,-0.921) (0.276,-0.906) (0.313,-0.889) (0.348,-0.870) (0.383,-0.849) (0.417,-0.827) (0.450,-0.803) (0.482,-0.778) (0.512,-0.751) (0.542,-0.724) (0.571,-0.695) (0.598,-0.665) (0.625,-0.635) (0.651,-0.603) (0.676,-0.571) (0.700,-0.539) (0.723,-0.505) (0.746,-0.472) (0.767,-0.438) (0.789,-0.403) (0.809,-0.368) (0.829,-0.333) (0.849,-0.297) (0.868,-0.262) (0.887,-0.226) (0.906,-0.190) (0.924,-0.154) (0.942,-0.117)};
\addplot[lineaE, flechaen=0.285, flechaen=0.715] coordinates {(-0.914,-0.098) (-0.887,-0.128) (-0.859,-0.158) (-0.832,-0.188) (-0.805,-0.218) (-0.777,-0.248) (-0.749,-0.277) (-0.720,-0.305) (-0.691,-0.334) (-0.662,-0.362) (-0.632,-0.389) (-0.602,-0.416) (-0.570,-0.442) (-0.539,-0.467) (-0.506,-0.491) (-0.473,-0.514) (-0.439,-0.536) (-0.404,-0.557) (-0.369,-0.576) (-0.333,-0.594) (-0.296,-0.611) (-0.258,-0.626) (-0.219,-0.638) (-0.180,-0.649) (-0.141,-0.658) (-0.101,-0.665) (-0.061,-0.669) (-0.020,-0.672) (0.020,-0.672) (0.061,-0.669) (0.101,-0.665) (0.141,-0.658) (0.180,-0.649) (0.219,-0.638) (0.258,-0.626) (0.296,-0.611) (0.333,-0.594) (0.369,-0.576) (0.404,-0.557) (0.439,-0.536) (0.473,-0.514) (0.506,-0.491) (0.539,-0.467) (0.570,-0.442) (0.602,-0.416) (0.632,-0.389) (0.662,-0.362) (0.691,-0.334) (0.720,-0.305) (0.749,-0.277) (0.777,-0.248) (0.805,-0.218) (0.832,-0.188) (0.859,-0.158) (0.887,-0.128) (0.914,-0.098)};
\addplot[lineaE, flechaen=0.255, flechaen=0.745] coordinates {(-0.882,-0.063) (-0.847,-0.082) (-0.812,-0.101) (-0.776,-0.119) (-0.741,-0.138) (-0.705,-0.156) (-0.669,-0.174) (-0.633,-0.191) (-0.597,-0.208) (-0.560,-0.225) (-0.523,-0.241) (-0.486,-0.257) (-0.449,-0.271) (-0.412,-0.285) (-0.374,-0.298) (-0.335,-0.311) (-0.297,-0.322) (-0.258,-0.332) (-0.219,-0.340) (-0.180,-0.348) (-0.140,-0.354) (-0.100,-0.358) (-0.060,-0.361) (-0.020,-0.363) (0.020,-0.363) (0.060,-0.361) (0.100,-0.358) (0.140,-0.354) (0.180,-0.348) (0.219,-0.340) (0.258,-0.332) (0.297,-0.322) (0.335,-0.311) (0.374,-0.298) (0.412,-0.285) (0.449,-0.271) (0.486,-0.257) (0.523,-0.241) (0.560,-0.225) (0.597,-0.208) (0.633,-0.191) (0.669,-0.174) (0.705,-0.156) (0.741,-0.138) (0.776,-0.119) (0.812,-0.101) (0.847,-0.082) (0.882,-0.063)};
\addplot[lineaE, flechaen=0.241, flechaen=0.759] coordinates {(-0.870,0.000) (0.870,0.000)};
\addplot[lineaE, flechaen=0.592] coordinates {(-1.130,0.000) (-2.600,0.000)};
\addplot[lineaE, flechaen=0.408] coordinates {(2.600,0.000) (1.130,0.000)};

  \fill[cargapos] (axis cs:-1,0) circle[radius=0.075];
  \fill[carganeg] (axis cs:1,0) circle[radius=0.075];
\end{axis}
\end{tikzpicture}
\caption{$+q$ y $-q$: $u_e$}
\end{subfigure}\hfill
\begin{subfigure}[b]{0.325\textwidth}\centering
\begin{tikzpicture}
\begin{axis}[mapa]
  \addplot graphics[xmin=-2.6, xmax=2.6, ymin=-1.8, ymax=1.8] {t5-ue-cruzado-\modotema};
  \draw[equipot] (axis cs:0,0) circle[radius=1];
  \fill[cargapos] (axis cs:-1,0) circle[radius=0.075];
  \fill[cargapos] (axis cs:1,0) circle[radius=0.075];
\end{axis}
\end{tikzpicture}
\caption{$+q$ y $+q$: $\eo\,\vE_1\cdot\vE_2$}
\end{subfigure}
\caption{La energía electrostática no cumple el principio de superposición. (a), (b) Densidad de energía
$u_e=\frac{\eo}{2}E^2$ (escala logarítmica, más brillante cuanto mayor) con las líneas de campo superpuestas: entre dos cargas
iguales el campo se anula en el punto medio y $u_e$ es pequeña; entre cargas opuestas los campos se refuerzan.
(c) Término cruzado $\eo\,\vE_1\cdot\vE_2$ para dos cargas iguales: es positivo (rojo) fuera de la circunferencia de
trazos, que tiene por diámetro el segmento entre las cargas, y negativo (azul) dentro. Su integral a todo el espacio
es exactamente la energía de interacción, $q^2/(4\pi\eo d)$ (\cref{ej:termino-cruzado}).}
\label{fig:superposicion}
\end{figure}


\begin{trampa}[Trampa habitual]
Si se duplica la carga de un sistema, su energía se multiplica por cuatro, no por dos; y la energía de dos sistemas juntos no es la suma
de sus energías por separado. Hay que calcular el campo total y elevarlo al cuadrado, o usar $\frac12\int\rho\phi\,\dd V$ con el potencial
total.
\end{trampa}
